\section{Model and Basic Quantities}\label{sec:foundations}
%==============================================================================

This section introduces only the minimum structure needed to state the deterministic converse foundation and its boundary threshold corollary: latent states, deterministic observations, auxiliary tags, ambiguity classes, and independent rate.

\subsection{Model assumptions and notation}\label{sec:assumptions}

We use the following standing assumptions.
\begin{enumerate}
\item \textbf{Finite latent state:} each fact takes values in a finite alphabet.
\item \textbf{Deterministic observations:} a system state exposes a deterministic observation transcript.
\item \textbf{Independent rate:} the independent rate counts the number of independently writable source locations for the same latent fact.
\item \textbf{Finite side information:} any auxiliary tag is drawn from a finite alphabet.
\end{enumerate}

These are the only ingredients needed for the deterministic zero-error theory developed below.

\subsection{Latent states, ambiguity, and independent rate}\label{sec:epistemic}

An \emph{encoding system} for a fact $F$ is a finite collection of locations $\{L_1,\dots,L_n\}$ whose current values jointly determine the observation transcript available to a resolver. For a system state $x$, the \emph{ambiguity class} is the set of latent values still compatible with that transcript; when the ambiguity class has size greater than one, the transcript alone does not identify a unique authoritative value. In this single-fact setting, latent-state ambiguity and value ambiguity coincide because the latent state is just the fact value. In the later multi-fact extension, ambiguity is instead over full latent tuples compatible with the partial observations. A \emph{fact} is an atomic semantic attribute of the represented object that can vary independently of other attributes, and a fact is \emph{structural} if the locations encoding it are fixed when the underlying object, declaration, or artifact is created.

The admissible edit set $E(C)$ is part of the primitive specification of an encoding system $C$. It records which edit events or edit sequences are considered reachable by the model. The paper does not derive $E(C)$ from a particular concurrency policy, protocol, or consensus mechanism; those are downstream instantiations. Independence and independent rate are therefore always defined \emph{relative to the chosen admissible edit model}.

\begin{definition}[Independent Location]\label{def:independent}
Let $E(C)$ denote the admissible edit set of encoding system $C$. Two locations encoding the same fact are \emph{independent relative to $E(C)$} if some admissible edit sequence in $E(C)$ can change one without forcing the other to match it, equivalently if some admissible edit sequence reaches a state in which the two locations disagree.
\end{definition}

\begin{definition}[Derived Location]\label{def:derived}
A location is \emph{derived} from a source location if updating the source determines the derived value without an additional manual edit.
\end{definition}

\begin{definition}[Degrees of Freedom]\label{def:dof}
The \emph{degrees of freedom} (DOF) is used as shorthand for the independent rate, defined as the number of pairwise independent source locations encoding $F$ under the specified admissible edit set $E(C)$. In the single-fact model, the remaining encoding locations are treated as derived views of those independent sources.
\end{definition}

\begin{definition}[Structural Integrity]\label{def:structural-integrity}
An encoding system has \emph{structural integrity} for fact $F$ if every reachable state is coherent, so no reachable system state presents mutually incompatible encodings of $F$.
\end{definition}

\begin{definition}[Integrity Violation]\label{def:integrity-violation}
A reachable \emph{integrity violation} is a reachable incoherent state for fact $F$.
\end{definition}

\leanmetapending{\LH{COH1}}
\begin{proposition}[Unit Independent Rate Guarantees Coherence]\label{thm:dof-one-coherence}
If the independent rate is $1$, then all reachable states are coherent.
\end{proposition}

\begin{proof}
At independent rate $1$, exactly one location is independent and all remaining locations are derived from it. Derived locations cannot diverge from their source, so disagreement is unreachable.
\end{proof}

\leanmetapending{\LH{COH2}}
\begin{proposition}[Independent Rate Above One Permits Incoherence]\label{thm:dof-gt-one-incoherence}
If the independent rate exceeds $1$, then incoherent states are reachable.
\end{proposition}

\begin{proof}
With at least two independent locations, admissible edits can assign different values to those locations, producing a reachable disagreement state.
\end{proof}

\subsection{Zero-Incoherence Threshold}\label{sec:capacity}

\begin{definition}[Zero-Incoherence Threshold]\label{def:coherence-capacity}
The \emph{zero-incoherence threshold} is the largest independent rate for which incoherent states are unreachable.
\end{definition}

\leanmetapending{\LH{RAT1}, \LHrng{CAP}{2}{3}}
\begin{corollary}[Threshold Corollary]\label{thm:coherence-capacity}
For deterministic multi-location encodings, the zero-incoherence threshold equals $1$.
\end{corollary}

\begin{proof}
\emph{Achievability.} Suppose $\mathrm{DOF}(C,F)=1$, and let $L_s$ be the unique independent location. By Definition~\ref{def:dof}, every other encoding location is non-independent and is therefore treated in this model as a derived view of $L_s$. By Definition~\ref{def:derived}, updating $L_s$ determines the value of each such derived location without an additional manual edit. Hence all encoding locations agree in every reachable state, so zero incoherence is achieved.

\emph{Converse.} Suppose $\mathrm{DOF}(C,F)=k>1$. By Definition~\ref{def:independent}, there exist at least two independent locations $L_1,L_2$ and an admissible edit sequence that can make them disagree. Equivalently, choose values $v_1\neq v_2$ and apply edits setting $L_1\leftarrow v_1$ and then $L_2\leftarrow v_2$. By independence, the second edit does not force $L_1$ to change. The resulting reachable state satisfies $\mathrm{value}(L_1)\neq \mathrm{value}(L_2)$ and is therefore incoherent.

Thus zero-error recovery is possible exactly at independent rate $1$, so the zero-incoherence threshold equals $1$.
\end{proof}

\leanmetapending{\LHrng{COH}{1}{2}}
\begin{corollary}[Exact Structural-Integrity Threshold]\label{cor:integrity-threshold}
A deterministic multi-location encoding has structural integrity if and only if its independent rate equals $1$.
\end{corollary}

\begin{proof}
If the independent rate is $1$, Proposition~\ref{thm:dof-one-coherence} shows that all reachable states are coherent, hence Definition~\ref{def:structural-integrity} holds. If the independent rate exceeds $1$, Proposition~\ref{thm:dof-gt-one-incoherence} gives a reachable incoherent state, so structural integrity fails.
\end{proof}

This threshold is an early architectural consequence of the deterministic model. It identifies the unique zero-error corner and, equivalently, the exact structural-integrity regime. Above that boundary, integrity violations are reachable by construction. The deeper theorems of the paper begin once one asks how much finite side information is needed whenever the transcript leaves a nontrivial ambiguity class.

\leanmetapending{\LH{CIA1}}
\begin{corollary}[Side-Information Lower Bound]\label{thm:side-info}
If a surviving ambiguity class has size $k$, then exact zero-error resolution requires at least $\log_2 k$ bits of side information.
\end{corollary}

\begin{proof}
Exact resolution of $k$ surviving alternatives requires at least $k$ distinguishable tags, hence at least $\log_2 k$ bits.
\end{proof}

The remainder of the paper sharpens this counting statement into a unified finite converse family, asks how that one-shot obstruction acquires non-clique topology under partial views, and then develops the corresponding asymptotic capacity and upper-bound theory.
